\documentclass[11pt]{article}

\usepackage[T1]{fontenc}
\usepackage[utf8]{inputenc}
\usepackage[margin=1in]{geometry}
\usepackage{xcolor}
\usepackage{enumitem}
\usepackage{hyperref}
\usepackage{listings}

\definecolor{codegray}{HTML}{F4F4F4}
\definecolor{linkblue}{HTML}{1F5AA6}

\hypersetup{
  colorlinks=true,
  linkcolor=linkblue,
  urlcolor=linkblue,
  pdftitle={IBM RAG and Agentic AI Lab Guide},
  pdfauthor={Djordje}
}

\lstset{
  backgroundcolor=\color{codegray},
  basicstyle=\ttfamily\small,
  breaklines=true,
  frame=single,
  showstringspaces=false
}

\title{IBM RAG and Agentic AI\\Lab-by-Lab Guide}
\author{Djordje}
\date{\today}

\begin{document}

\maketitle

\begin{abstract}
This document explains each lab in the IBM RAG and Agentic AI showcase.
It focuses on the problem each lab solves, how data moves through the
pipeline, why the main engineering decisions matter, and how each result
supports the eventual retrieval-augmented and agentic system. The guide is
updated as new labs are added.
\end{abstract}

\tableofcontents
\newpage

\section{Lab 01: Structured Restaurant Extraction}

\subsection{Goal}

The first lab converts free-form restaurant descriptions into consistent JSON.
Human-readable prose contains useful details, but fields such as location,
rating, signature dishes, atmosphere, and shortcomings cannot be filtered or
indexed reliably until they have a stable structure.

\subsection{Input and output}

The input is the \emph{California Culinary Map}, a text document containing 210
restaurant descriptions. The output is a JSON array. Each item contains:

\begin{itemize}[nosep]
  \item name, location, restaurant type, and food style;
  \item optional rating, price range, and vibe;
  \item lists of signature items and shortcomings;
  \item a required environment description; and
  \item a deterministic \texttt{itemId}.
\end{itemize}

\subsection{How the lab works}

\begin{enumerate}
  \item Download the text dataset and split it at blank lines.
  \item Remove the title block, leaving one paragraph per restaurant.
  \item Build a one-shot prompt containing the required fields, explicit rules,
        and one representative input/output example.
  \item Send the prompt to IBM Granite through watsonx.ai.
  \item Validate the returned JSON with a Pydantic \texttt{Restaurant} model.
  \item If validation fails, send the candidate and exact validation errors to
        a repair prompt.
  \item Stop after the configured number of repair attempts rather than
        retrying forever.
  \item Add stable item IDs and write the validated collection to disk.
\end{enumerate}

\subsection{Why validation matters}

An LLM can return syntactically valid JSON that still has missing fields or
incorrect types. Pydantic treats model output as untrusted input. Only records
that satisfy the schema reach the final dataset. Validation errors also provide
specific repair instructions, which are more useful than asking the model to
``try again'' without context.

\subsection{Important engineering decisions}

\begin{itemize}
  \item Model access is injected as a callable, so tests can use fake responses
        without network access or inference cost.
  \item Repair attempts are bounded to prevent infinite loops.
  \item Downloaded data and generated outputs are excluded from Git.
  \item Credentials are read from environment variables and are never stored
        in source code.
\end{itemize}

\subsection{Run Lab 01}

\begin{lstlisting}[language=bash]
pip install -e ".[dev,labs]"
export WATSONX_APIKEY="..."
export WATSONX_PROJECT_ID="..."
restaurant-extract --limit 3
\end{lstlisting}

\section{Lab 02: Multimodal Food-Data Augmentation}

\subsection{Goal}

The second lab converts visual information into text that can later be embedded
and retrieved. It augments recipes with descriptions of finished dishes and
augments restaurant reviews with captions of attached images.

\subsection{Input and output}

The recipe input contains 109 structured recipe records and a 215 MB archive of
matching images. Recipe ID \(n\) maps to
\texttt{synthetic\_recipe\_images/recipe\(n\).png}. Each output recipe gains an
\texttt{image\_description} field.

The review input contains 10 synthetic user reviews. Its \texttt{images} field
is a string representation of a URL list, so the lab parses it safely with
\texttt{ast.literal\_eval}. Each output review gains an
\texttt{image\_captions} list.

\subsection{How the recipe pipeline works}

\begin{enumerate}
  \item Download the recipe JSON and image archive when they are not cached.
  \item Extract the archive only after checking every member for path traversal
        and symbolic links.
  \item Resolve the image path from the recipe's integer ID.
  \item Build a caption prompt using the dish name.
  \item Encode the image as a base64 data URL with the correct MIME type.
  \item Ask the vision-language model for a concise caption covering visible
        ingredients, colors, plating, and style.
  \item Add the caption as \texttt{image\_description}.
\end{enumerate}

\subsection{How the review pipeline works}

\begin{enumerate}
  \item Parse each review's stringified image URL list.
  \item Download every image with bounded exponential retry.
  \item Read the review from the live dataset's \texttt{text} field.
  \item Give both the image and review text to the vision-language model.
  \item Ask for details that explain the reviewer's preference or experience.
  \item Store all successful captions in \texttt{image\_captions}; skip an
        image if its download still fails after all retries.
\end{enumerate}

\subsection{Why contextual captioning matters}

Pixels alone may show a plate, drink, or dining room without explaining why it
was meaningful to the reviewer. Supplying review text gives the model grounded
context. The resulting caption can connect visual evidence---for example,
greenhouse decor or elaborate plating---to the written experience.

\subsection{Reliability and cost controls}

\begin{itemize}
  \item One model client is reused instead of being reconstructed per image.
  \item ZIP extraction rejects unsafe archive members.
  \item Image MIME types are detected rather than assumed to be JPEG.
  \item Network retry is exponential, bounded, and limited to request failures.
  \item The example processes three records by default. Setting
        \texttt{LAB02\_LIMIT=0} explicitly enables the complete batch.
\end{itemize}

\subsection{Run Lab 02}

\begin{lstlisting}[language=bash]
pip install -e ".[labs]"
export WATSONX_APIKEY="..."
export WATSONX_PROJECT_ID="..."

# Open this file as Python notebook cells:
examples/02_multimodal_food_data_augmentation.py
\end{lstlisting}

\section{Lab 03: Safe Restaurant Database}

\subsection{Goal}

The third lab turns Lab 01's structured JSON into an interactive terminal
application. A user can browse restaurants, view details, add a restaurant from
natural-language prose, edit an existing record, or delete a record. The main
challenge is not the menu itself; it is preserving schema validity and data
safety while accepting human input.

\subsection{How the application is organized}

The lab separates three responsibilities:

\begin{itemize}
  \item \textbf{Storage} loads validated records, creates backups, and replaces
        the JSON file atomically.
  \item \textbf{Database operations} implement add, update, delete, and stable
        ID allocation without depending on terminal input.
  \item \textbf{Terminal interface} displays the menu, parses user input,
        confirms writes, and reports errors without crashing.
\end{itemize}

This separation makes the important behavior testable without calling a model
or manually interacting with the menu.

\subsection{Browse and view operations}

Browse and view are read-only. The application loads every JSON object as a
Pydantic \texttt{RestaurantRecord}. If an existing record is malformed, loading
fails visibly rather than allowing corrupted data to propagate. Read-only use
does not initialize watsonx.ai and does not require cloud credentials.

\subsection{Adding a restaurant}

\begin{enumerate}
  \item The user enters write mode and types \texttt{yes} to confirm.
  \item The user supplies a natural-language restaurant description.
  \item The application lazily initializes the Lab 01 Granite extractor.
  \item Granite returns structured JSON, which the Lab 01 schema validates and
        repairs with a bounded retry loop when necessary.
  \item The database chooses one plus the highest existing \texttt{itemId}.
  \item The current JSON is backed up and the expanded collection is written
        through a temporary file.
\end{enumerate}

Using the highest ID avoids a subtle collision. An ID based only on list length
can reuse an existing value after an earlier record has been deleted.

\subsection{Editing a restaurant}

Terminal input begins as text, but database fields do not all have string types.
The lab converts ratings to floating-point values, price ranges to integers,
\texttt{null} to \texttt{None}, and comma-separated or JSON-formatted list
fields to lists of strings. It then validates the \emph{complete} updated record
before saving. The stable \texttt{itemId} is not editable.

\subsection{Deleting a restaurant}

Deletion uses a one-based restaurant number for the user and checks that the
corresponding zero-based index exists. The write-mode confirmation occurs before
the mutation. A cancelled or invalid deletion leaves both the in-memory data
and file unchanged.

\subsection{Backup and atomic replacement}

Immediately before a successful write, the current JSON file is copied to a
\texttt{.bak} path. New JSON is first written to a temporary file and then moved
over the database path with an atomic replacement. This prevents a partially
written main file if serialization or writing fails midway.

\subsection{Run Lab 03}

\begin{lstlisting}[language=bash]
# First create a small structured dataset with Lab 01.
restaurant-extract --limit 5

# Then open the interactive database.
restaurant-db \
  --file data/processed/structured_restaurant_data.json
\end{lstlisting}

\subsection{Offline testing}

Tests use temporary directories and a fake restaurant extractor. They verify
backup behavior, stable ID generation, typed edits, forbidden field changes,
deletion, and cancellation. No test consumes inference credits.

\section{How the Labs Connect}

Lab 01 creates clean, validated restaurant records. Lab 02 enriches textual
records with information extracted from images. Lab 03 adds a safe human
interface for maintaining the structured restaurant collection. Together they
prepare a corpus with stable identifiers, structured metadata,
natural-language descriptions, visual captions, and controlled updates.

The next retrieval-augmented generation lab can chunk and embed these fields,
store them in a vector index, retrieve relevant evidence for a user's question,
and produce grounded recommendations. Later agentic labs can add planning and
tools around that retrieval layer.

\section{Repository and Reproducibility}

The complete source is available at:

\begin{center}
\url{https://github.com/Djordje3002/ibm-rag-agentic-ai-showcase}
\end{center}

Unit tests use fake model functions, so validation, prompt flow, path handling,
and augmentation logic can be checked without cloud credentials. Actual model
inference remains an explicit, credentialed step.

\end{document}
